%% MOTIVAȚIE
\chapter*{Motivație}
\addcontentsline{toc}{chapter}{Motivație}

Procesul care a dus la dezvoltarea acestei aplicații nu a început dintr-o
intenție pur tehnică, ci dintr-o observație persistentă a problemelor întâlnite
în practica de zi cu zi a mediului universitar. Pe parcursul anilor de
facultate, m-am angajat activ în diverse activități didactice și
administrative, ceea ce mi-a oferit perspectiva atât a unui student, cât și a
unui observator al proceselor instituționale. În acest context, platforma
Moodle a fost instrumental importantă: a servit ca spațiu central pentru
distribuirea materialelor de curs, comunicarea temelor și sarcinilor de lucru,
precum și pentru urmărirea activităților și progresului didactic. Moodle
rămâne, fără îndoială, o soluție utilă și solid construită pentru domeniul
e-learning, cu o comunitate de dezvoltare activă și o bază de utilizatori
impresionantă la nivel mondial.

Totuși, în interacțiunea de la un semestru la altul cu această platformă și cu
procesele academice din jur, am observat că, deși Moodle acoperă bine aspectele
de curs tradițional, în practica de zi cu zi a activităților de laborator și a
administrării operaționale, situațiile reale ajung să genereze fragmente
necontrolate de flux. Unele informații critice rămân înscrise în platformă,
altele migrează treptat în tabele locale gestionate de cadre didactice, în
fișiere Excel care circulă prin email sau sunt stocate pe calculatoare
personale, iar altele încă se evidențiază manual în caiete de prezență. Această
dispersie nu este o eroare a vreunui sistem individual, ci mai degrabă un
decalaj sistemic între ceea ce o platformă generalistă poate oferi și ceea ce
practicienii universității au nevoie în momentul concret.

Această fragmentare generează o fricțiune administrativă care nu este nici
spectaculoasă, nici catastrofală, dar nici neglijabilă. O sarcină care ar putea
să fie simplă — actualizarea unei liste de prezență pe o sesiune de laborator
sau centralizarea studenților unui grup pentru o activitate nouă — ajunge să
implice mai mulți pași manuali, verificări repetate și sincronizări între surse
disparate. Pe măsură ce volumul de date și de tranzacții crește pe parcursul
semestrului, costul acumulat al acestor operații devine evident: apar
întârzieri în procesarea informațiilor, apar diferențe semnificative între
versiunile de date stocate în locuri diferite, iar luarea de decizii devine mai
laborioasă din cauza lipsei unei imagini unice și actualizate asupra situației
reale din teren. Din perspectiva utilizatorilor finali, experiența se
fragmentează și devine inconsistentă: un student trebuie să caute informații în
mai multe locuri, un profesor se află în situația de a repeta aceeași operație
în mai multe sisteme în paralel, iar un administrator dedică o proporție
disproporționată a timpului său gestionării cazurilor speciale și excepțiilor,
în loc să poată construi și îmbunătăți procese standardizate și scalabile.

Motivația principală a acestei lucrări a fost, deci, construirea unei soluții
specific adaptate acestui context local și acestor necesități reale, o
aplicație care să completeze ceea ce lipsește din practica curentă și care să
reducă substanțial dependența de fluxuri paralele, neintegrate și sincronizate
manual. Ideea de bază nu a fost niciodată aceea de a înlocui o platformă
consacrată și bine stabilită, ci mai degrabă de a proiecta o aplicație
verticalizată, focalizată exclusiv pe managementul fluxurilor studențești în
sens operațional, în particular în acele zone în care viteza de operare,
claritatea, coerența datelor și trasabilitatea auditabilă sunt esențiale pentru
buna desfășurare a activităților.

Academo a fost concepută din start ca o aplicație web în care operațiile
frecvente, repetitive, să devină naturale și intuitive de executat:
autentificarea utilizatorilor cu diferențiere clară de roluri, administrarea și
organizarea studenților în grupe funcționale, organizarea și programarea
cursurilor și sesiunilor de laborator, marcarea și urmărirea prezenței în mod
fluid și direct la locul de activitate, importul rapid de date din fișiere
Excel și accesul rapid la rapoarte și statistici. Din perspectiva experienței
utilizatorului, obiectivul central a fost ca interfața să fie previzibilă și
coerentă în toți pașii, suficient de flexibilă încât să se adapteze la ritmul
și la particularitățile reale ale activităților din cadrul facultății, dar și
suficient de structurată pentru a proteja utilizatorii de erori și a-i ghida
spre acțiunile corecte.

Pe lângă componenta practică și imediat utilă, această motivație are și o
dimensiune profundă de formare și dezvoltare profesională personală.
Dezvoltarea acestei aplicații a reprezentat un cadru ideal și un pretext
potrivit pentru a aprofunda și a consolida cunoștințele în tehnologiile moderne
de frontend web, în special în framework-ul Angular și ecosistemul său.
Totodată, a oferit oportunitatea de a aplica și testa în practică principii
fundamentale de arhitectură software și design în contextul unui proiect
cuprinzător și real, nu doar într-un prototip izolat sau într-un exercițiu
academic. În sens mai larg, lucrarea combină oportun nevoia reală și tangibilă
din mediul universitar cu aspirația personală de a construi o soluție robust
proiectată, extensibilă și ușor de menținut și evoluat în viitor.

\newpage

\chapter*{Introducere}
\addcontentsline{toc}{chapter}{Introducere}
Digitalizarea administrației academice în instituțiile de învățământ superior nu se reduce la un simplu transfer al formularelor tradiționale într-un mediu electronic, ci presupune o regândire și o reproiectare fundamentală a fluxurilor de lucru, a modului în care informația circulă între actorii implicați și a felului în care operațiile repetitive pot fi reduse la minimum prin automatizare inteligentă și prin design chibzuit. În forma lor clasică și inerțială, procesele administrative precum înscrierea și organizarea studenților în grupe de lucru, pregătirea și programarea sesiunilor de laborator, notarea prezenței participanților sau centralizarea și raportarea situațiilor academice, au fost tradițional distribuite pe mai multe canale și instrumente disparate. Atunci când lipsește un sistem unificat care să integreze aceste procese și să le asigure o bază comună de date, consistența și corectitudinea informațiilor depind în mare măsură de disciplina și de buna pregătire a indivizilor, mai degrabă decât de regulile tehnice impuse de o platformă standardizată și automatizată.

Lucrarea de față propune proiectarea și implementarea tehnică a unei aplicații
web dedicate satisfacerii acestor nevoi concrete și operaționale, o arhitectură
client-server în care componentele sunt bine delimitate și responsabilitățile
sunt clar definite. Frontend-ul, interfața cu care utilizatorul interacționează
zilnic, este realizat cu ajutorul framework-ului Angular, o tehnologie modernă
pentru construirea interfețelor web interactive și responsive. Backend-ul,
stratul în care rezidă logica de business și deciziile critice ale aplicației,
expune servicii prin intermediul unui API RESTful implementat în limbajul PHP,
care este bine cunoscut, ușor de integrat și larg suportat în infrastructura
academică. Persistența datelor — stocarea și integritatea pe termen lung a
informațiilor — este asigurată prin intermediul unei baze de date relaționale
MySQL, o tehnologie stabilă, scalabilă și cu suport solid pe platformele LAMP
tradiționale. Alegerea acestei arhitecturi, cu separarea clară între frontend,
backend și baza de date, a fost motivată de dorința de a realiza o decuplare
riguroasă a responsabilităților: interfața utilizatorului și experiența
acestuia sunt gestionate exclusiv de frontend, logica complexă de business și
validările critice sunt centralizate în backend, iar integritatea structurală
și consistența semantică a datelor sunt garantate de stratul bazei de date.

Un accent deosebit al acestei lucrări este pus pe aspectele de frontend și pe
framework-ul Angular, deoarece în această componentă se regăsesc și se
materializează majoritatea deciziilor de arhitectură care influențează direct
experiența utilizatorului și mentenabilitatea pe termen lung a soluției. Sunt
discutate și analizate teme importante: structurarea rațională a componentelor
și modularitatea ierarhică, organizarea logică a modulelor funcționale,
mecanismele de rutare și protecția rutelor de acces, comunicarea între
componente prin intermediul serviciilor și injecția de dependențe, lucrul
reactiv și declarativ cu fluxurile de date observabile (RxJS), validarea
avansată și reactivă a formularelor prin Reactive Forms și integrarea
comunicării HTTP cu interceptori pentru gestionarea unificată a autentificării
și a erorilor. Aceste elemente nu sunt prezentate abstract, ca o simplă
enumerare de concepte sau caracteristici, ci sunt analizate și exemplificate ca
piese care se leagă armonic în cadrul unui flux funcțional complet și coerent.

Din punct de vedere funcțional și al cazurilor de utilizare, aplicația
realizată acoperă setul de procese considerate esențiale pentru operarea unei
facultăți sau a unui departament: mecanismele de autentificare și autorizare pe
bază de roluri diferențiate, administrarea și managementul studenților și al
grupelor de lucru, gestionarea cursurilor și a sesiunilor de laborator
asociate, marcarea și urmărirea prezenței pe sesiuni, importul și actualizarea
în masă a datelor din fișiere Excel, și vizualizarea de rapoarte și statistici
sintetice. Fiecare modul funcțional al aplicației este analizat atât din
perspectiva tehnică și de implementare, cât și din perspectiva utilizării
practice și a impactului asupra experienței utilizatorului, pentru a evidenția
și explica felul în care alegerile tehnice de implementare influențează viteza
și fluiditatea operării, acuratețea și consistența datelor, precum și
claritatea și intuitivitatea interfeței.

Din perspectiva metodologică și a procesului de dezvoltare, realizarea acestei
aplicații a urmat un traseu ponderat și incremental, care a inclus etape
succesive: definirea riguroasă a cerințelor și a constrângerilor, modelarea
riguroasă a datelor și a relațiilor dintre entități, proiectarea și
documentarea arhitecturii sistemului, implementarea progresivă a modulelor și a
componentelor, validarea soluției prin scenarii reale de utilizare și prin
feedback-ul utilizatorilor pilot, precum și ajustări iterative pe baza
observațiilor și a lecțiilor învățate. Această abordare incrementală și
iterativă a permis rafinarea continuă și evoluția soluției, menținând un
echilibru dinamic între cerințele și constrângerile tehnice, pe de o parte, și
utilitatea practică și relevanța reală a soluției în contextul concret
universitar, pe de altă parte.

Structura lucrării este organizată în patru capitole principale, fiecare
servind un scop bine definit în construirea unei înțelegeri de ansamblu a
temei. Capitolul întâi prezintă și analizează tehnologiile și framework-urile
utilizate în dezvoltarea acestei aplicații, cu un accent extins și detaliat pe
framework-ul Angular și pe mecanismele sale arhitecturale care sunt relevante
și aplicabile în contextul acestui proiect. Capitolul al doilea descrie pe
îndelete arhitectura sistemului, modelul relațional al datelor și furnizează
diagramele și explicațiile care clarifică relațiile și fluxurile de date dintre
componentele sistemului. Capitolul al treilea oferă o prezentare narativă și
practică a aplicației prin fluxuri concrete de utilizare reală, de la momentul
inițial al autentificării utilizatorului, trecând prin importul de date din
fișiere Excel, până la procesele de marcare a prezenței și accesul la rapoarte.
Capitolul al patrulea, cel final, sintetizează rezultatele efectiv obținute
prin implementare, identifică limitările și punctele forte ale sistemului și
propune direcții viabile de dezvoltare și îmbunătățire pe viitor. Prin această
structură coerentă și progresivă, lucrarea nu doar demonstrează că aplicația
este funcțională și operațională, ci explică și justifică de ce soluția
construită este potrivită și bine adaptată problemei adresate, cum poate fi
extinsă și adaptată la contexte și cerințe în evoluție, și în ce mod aplicația
poate susține efectiv și durabil procesele academice în contexte reale pe
termen lung. 
\newpage
